\color{uniblue}\section[Основные функции]{Основные функции}\label{sec:funcs}
\color{black}

\needspace{4\baselineskip}
\color{uniblue}{\subsection{Общие сведения}}
\color{black}

\begin{enumerate}[label=\arabic{section}.\arabic{subsection}.\arabic*, labelsep=4pt, leftmargin=0pt, itemindent=57pt]
    
\item
При вводе в работу устройство постоянно отслеживает состояние подключенных к нему аналоговых и дискретных сигналов. Из значений, полученных АЦП, при помощи расчетов и, при необходимости, цифровой фильтрации выделяются среднеквадратичные и/или мгновенные величины (фазных токов, гармонических составляющих токов, фазных и междуфазных напряжений), необходимые для работы функций в составе устройства.
\item
Оперативный вывод всех функций в составе устройства из работы происходит при появлении активного сигнала на входе <<Вывод терминала>>. 
\item
Значения задаваемых параметров для токов в таблицах параметров настройки функций приведены по отношению к номинальному значению тока аналогового входа и указываются в относительных единицах (о.е). 
\item
Общая схема взаимосвязей между описанными в данном разделе функциями, входными и выходными сигналами устройства приведена в приложении \ref{app:fsu}.
Сводный перечень сигналов, формируемых устройством в процессе работы, приведен в таблице \ref{appA:tbl1} приложения \ref{app:signals}.

\end{enumerate}

